\documentclass[10pt,a4paper]{amsart}
\usepackage{amsmath,amssymb,amsthm,graphicx}
\usepackage{verbatim}
\newenvironment{reference}{\comment}{\endcomment}
%\newenvironment{reference}{}{}
\newenvironment{slogan}{\comment}{\endcomment}
\newenvironment{history}{\comment}{\endcomment}

% For commutative diagrams we use Xy-pic
\usepackage[all]{xy}

% We use 2cell for 2-commutative diagrams.
\xyoption{2cell}
% Theorem environments.
%
\theoremstyle{plain}
\newtheorem{theorem}[subsection]{Theorem}
\newtheorem{proposition}[subsection]{Proposition}
\newtheorem{lemma}[subsection]{Lemma}

\theoremstyle{definition}
\newtheorem{definition}[subsection]{Definition}
\newtheorem{example}[subsection]{Example}
\newtheorem{exercise}[subsection]{Exercise}
\newtheorem{situation}[subsection]{Situation}

\theoremstyle{remark}
\newtheorem{remark}[subsection]{Remark}
\newtheorem{remarks}[subsection]{Remarks}

\numberwithin{equation}{subsection}

% Macros
%
\def\lim{\mathop{\mathrm{lim}}\nolimits}
\def\colim{\mathop{\mathrm{colim}}\nolimits}
\def\Spec{\mathop{\mathrm{Spec}}}
\def\Hom{\mathop{\mathrm{Hom}}\nolimits}
\def\Ext{\mathop{\mathrm{Ext}}\nolimits}
\def\SheafHom{\mathop{\mathcal{H}\!\mathit{om}}\nolimits}
\def\SheafExt{\mathop{\mathcal{E}\!\mathit{xt}}\nolimits}
\def\Sch{\mathit{Sch}}
\def\Mor{\mathop{\mathrm{Mor}}\nolimits}
\def\Ob{\mathop{\mathrm{Ob}}\nolimits}
\def\Sh{\mathop{\mathit{Sh}}\nolimits}
\def\NL{\mathop{N\!L}\nolimits}
\def\CH{\mathop{\mathrm{CH}}\nolimits}
\def\proetale{{pro\text{-}\acute{e}tale}}
\def\etale{{\acute{e}tale}}
\def\QCoh{\mathit{QCoh}}
\def\Ker{\mathop{\mathrm{Ker}}}
\def\Im{\mathop{\mathrm{Im}}}
\def\Coker{\mathop{\mathrm{Coker}}}
\def\Coim{\mathop{\mathrm{Coim}}}

% Boxtimes
%
\DeclareMathSymbol{\boxtimes}{\mathbin}{AMSa}{"02}

%
% Macros for moduli stacks/spaces
%
\def\QCohstack{\mathcal{QC}\!\mathit{oh}}
\def\Cohstack{\mathcal{C}\!\mathit{oh}}
\def\Spacesstack{\mathcal{S}\!\mathit{paces}}
\def\Quotfunctor{\mathrm{Quot}}
\def\Hilbfunctor{\mathrm{Hilb}}
\def\Curvesstack{\mathcal{C}\!\mathit{urves}}
\def\Polarizedstack{\mathcal{P}\!\mathit{olarized}}
\def\Complexesstack{\mathcal{C}\!\mathit{omplexes}}
% \Pic is the operator that assigns to X its picard group, usage \Pic(X)
% \Picardstack_{X/B} denotes the Picard stack of X over B
% \Picardfunctor_{X/B} denotes the Picard functor of X over B
\def\Pic{\mathop{\mathrm{Pic}}\nolimits}
\def\Picardstack{\mathcal{P}\!\mathit{ic}}
\def\Picardfunctor{\mathrm{Pic}}
\def\Deformationcategory{\mathcal{D}\!\mathit{ef}}
\usepackage{cleveref}
\crefname{originalequation}{original equation}{original equations}
\Crefname{originalequation}{Original equation}{Original equations}
\makeatletter
\newlabel{equation-bijective}{{06IM}{source}{}{sourcepointer-1}{}}
\newlabel{equation-bijective@cref}{{[originalequation][0][]06IM}{[1][0][]source}{}{}{}}
\makeatother
\crefname{originalreference}{original reference}{original references}
\Crefname{originalreference}{Original reference}{Original references}
\makeatletter
\newlabel{example-tangent-space-prorepresentable-functor}{{06IC}{source}{}{sourcepointer-2}{}}
\newlabel{example-tangent-space-prorepresentable-functor@cref}{{[originalreference][0][]06IC}{[1][0][]source}{}{}{}}
\makeatother
\crefname{originalreference}{original reference}{original references}
\Crefname{originalreference}{Original reference}{Original references}
\makeatletter
\newlabel{lemma-versal-object-S1}{{06SY}{source}{}{sourcepointer-3}{}}
\newlabel{lemma-versal-object-S1@cref}{{[originalreference][0][]06SY}{[1][0][]source}{}{}{}}
\makeatother
\title{A predeformation category with a versal object satisfying the bijectivity condition satisfies (S1), (S2), and has finite dimensional tangent space}
\author{The Stacks Project Authors}
\date{stacks-project:formal-defos.tex:local-snapshot-20261012; GNU FDL 1.2 or later}
\begin{document}
\maketitle
\noindent\textit{Editorial note.} All delivered mathematical text is the author's own, copied from the pinned Stacks Project chapter file formal-defos.tex; no statement or proof step was rewritten, reordered or removed.\par\medskip
\noindent\textit{Editorial note.} The counted proof is the author's complete argument: the importation of (S1) and the first half of (S2), the proof of the second half of (S2) through the pushforwards of the versal object and the bijectivity condition, and the finiteness of the tangent space from the prorepresentable computation.\par\medskip
\noindent\textit{Editorial note.} The label lemma-miniversal-object-existence-1 names the counted statement, and the proof environment immediately following it in the source is delivered verbatim as its proof.\par\medskip
\noindent\textit{Editorial note.} The counted statement is carried by the author's own lemma environment in the chapter source.\par\medskip
\noindent\textit{Editorial note.} Pre-open gate: python operations/agent-authoring/tools/precard\_gate.py --source formal-defos.tex --label lemma-miniversal-object-existence-1 --start 4664 --end 4729 --local-id 28026 returned RESULT=OK (L28026 free, occupancy FREE, public interval clear). Candidate taken from sources/stacks-formal-defos/candidates-halfA.json.\par\medskip
\noindent\textit{Editorial note.} The statement and proof are copied from the pinned chapter file; the cited inputs are kept as external pointers with their canonical Stacks tags.\par\medskip
\noindent\textit{Allowed standard prerequisites (editorial).}\begin{itemize}
\item (S1) for a predeformation category with a versal object. Same chapter, the lemma deriving condition (S1) from the existence of a versal formal object satisfying the bijectivity condition. Supplies the first assertion and the first half of (S2).
\item Diagrams for the second half of (S2) and the morphisms from the versal object. Same chapter, the construction taking a diagram over the square with k[epsilon] and producing morphisms from the versal formal object into its objects. Supplies the morphisms b and b' from xi into y and y' with the commuting square.
\item Prorepresentability and finiteness of the tangent space. Same chapter, the example computing the tangent space of a prorepresentable functor, which is finite dimensional. Supplies the finiteness of the dimension of the tangent space of the predeformation category.
\end{itemize}\medskip
\ifdefined\hypertarget\hypertarget{sourcepointer-1}{}\fi
\noindent\textit{Source reference.} Original equation 06IM: sources/stacks-formal-defos/pinned-original-source.tex, source lines 4617--4617. Prerequisite pointer invoked at proof line(s) 4668; it names a standard fact of the same chapter that the counted argument uses, and every step of the counted proof built on it is retained verbatim in the delivered body.\par\medskip
\ifdefined\hypertarget\hypertarget{sourcepointer-2}{}\fi
\noindent\textit{Source reference.} Original reference 06IC: sources/stacks-formal-defos/pinned-original-source.tex, source lines 3545--3545. Prerequisite pointer invoked at proof line(s) 4725; it names a standard fact of the same chapter that the counted argument uses, and every step of the counted proof built on it is retained verbatim in the delivered body.\par\medskip
\ifdefined\hypertarget\hypertarget{sourcepointer-3}{}\fi
\noindent\textit{Source reference.} Original reference 06SY: sources/stacks-formal-defos/pinned-original-source.tex, source lines 3894--3894. Prerequisite pointer invoked at proof line(s) 4678; it names a standard fact of the same chapter that the counted argument uses, and every step of the counted proof built on it is retained verbatim in the delivered body.\par\medskip
\begin{lemma}
\label{lemma-miniversal-object-existence-1}
Let $\mathcal{F}$ be a predeformation category.
Let $\xi$ be a versal formal object of $\mathcal{F}$ such that
(\ref{equation-bijective}) holds. Then
\begin{enumerate}
\item $\mathcal{F}$ satisfies (S1).
\item $\mathcal{F}$ satisfies (S2).
\item $\dim_k T\mathcal{F}$ is finite.
\end{enumerate}
\end{lemma}

\begin{proof}
Condition (S1) holds by
Lemma \ref{lemma-versal-object-S1}.
The first part of (S2) holds since (S1) holds. Let
$$
\vcenter{
\xymatrix{
y \ar[r]_c \ar[d]_a & x_\epsilon \ar[d]^e \\
x \ar[r]^d          & x_0
}
}
\quad\text{and}\quad
\vcenter{
\xymatrix{
y' \ar[r]_{c'} \ar[d]_{a'} & x_\epsilon \ar[d]^e \\
x \ar[r]^d                 & x_0
}
}
\quad\text{lying over}\quad
\vcenter{
\xymatrix{
A \times_k k[\epsilon] \ar[r] \ar[d] & k[\epsilon] \ar[d] \\
A  \ar[r] & k
}
}
$$
be diagrams as in the second part of (S2). As above we can find
morphisms $b : \xi \to y$ and $b' : \xi \to y'$ such that
$$
\xymatrix{
\xi \ar[r]^{b'} \ar[d]_b          & y' \ar[d]^{a'} \\
y \ar[r]^{a} & x
}
$$
commutes. Let $p : \mathcal{F} \to \mathcal{C}_\Lambda$ denote the
structure morphism. Say $\widehat{p}(\xi) = R$, i.e., $\xi$ lies over
$R \in \Ob(\widehat{\mathcal{C}}_\Lambda)$. We see that the
pushforward of $\xi$ via $p(c) \circ p(b)$ is $x_\epsilon$
and that the pushforward of $\xi$ via $p(c') \circ p(b')$ is $x_\epsilon$.
Since $\xi$ satisfies (\ref{equation-bijective}), we see that
$p(c) \circ p(b) = p(c') \circ p(b')$
as maps $R \to k[\epsilon]$. Hence $p(b) = p(b')$ as maps from
$R \to A \times_k k[\epsilon]$. Thus we see that $y$ and $y'$ are
isomorphic to the pushforward of $\xi$ along this map and we get
a unique morphism $y \to y'$ over $A \times_k k[\epsilon]$
compatible with $b$ and $b'$ as desired.

\medskip\noindent
Finally, by
Example \ref{example-tangent-space-prorepresentable-functor}
we see
$\dim_k T\mathcal{F} = \dim_k T\underline{R}|_{\mathcal{C}_\Lambda}$
is finite.
\end{proof}

\end{document}
